% Source: https://artofproblemsolving.com/community/c3h1976306_triangle_inequality
% Stable identity: aops-topic-1976306-q01
Let the third side be $x$. By the triangle inequality, each side is less than the sum of the other two:
\[
x<20+19=39,\qquad 20<19+x,\qquad 19<20+x .
\]
The second inequality gives $x>1$ and the third holds automatically. So $1<x<39$, and the integer values are $x=2,3,\dots,38$.

\begin{asy}
usepackage("amsmath");
size(4.5cm);
defaultpen(fontsize(10pt));
pair B=(0,0),C=(20,0),A=(7.5,15.0);
draw(A--B--C--cycle,linewidth(0.9));
label("$20$",(B+C)/2,S);label("$19$",(A+C)/2,NE);label("$x$",(A+B)/2,NW);
dot(A);dot(B);dot(C);
label("\mbox{}",(min(currentpicture,user=true)+max(currentpicture,user=true))/2,black); // keeps the text after the figure black when LaTeX typesets the labels inline
\end{asy}

Each of these lengths gives exactly one triangle (up to congruence), so there are $38-2+1=\boxed{37}$ triangles.

(The bound $11<x<39$ in the post is wrong: the lower bound is $20-19=1$, not $20-9$.)
